%!TEX root = ../../csuthesis_main.tex
\chapter{绪论}
\section{研究背景与意义}
\subsection{研究背景}
随着水域无人载具（包括无人水面艇 USV 与水下航行器 AUV/ROV）在海洋勘探、环境监测与水域安防等领域的应用不断深入，其自主作业能力正从“单任务遥控”向“多任务自主规划与应急处置”演进。在实际任务中，水域载具需要具备航线规划、覆盖搜索、避障、动态重规划与返航策略等完整的任务规划能力。然而，复杂多变的水下环境往往伴随着诸多不确定性因素，推进器故障、定位漂移、传感器异常与通信中断等典型异常工况是水域载具在执行任务过程中必须面对的严峻挑战。

当前，多数任务规划算法的验证仍停留在理想仿真或单场景试验阶段。算法在无故障条件下“跑通”并不等同于在异常叠加条件下仍能维持任务连续性或安全返航。与此同时，推进器失效、传感器漂移、通信链路中断等故障场景在真实水域中不仅难以低成本复现，且往往伴随着极高的安全风险。因此，构建高保真的虚拟仿真环境成为了解决这一痛点的有效途径。近年来，以虚幻引擎（Unreal Engine）为代表的高保真渲染技术，为构建逼真的水域环境与传感器模型提供了底层支撑，也为开展系统级的任务规划验证与故障注入提供了理想的试验场。

\subsubsection{研究意义}

针对上述问题，构建一个能够同时支撑“任务规划全流程验证”与“多类型故障注入”的仿真验证系统，具有明确的工程与学术意义。

在\textbf{理论意义}方面，本课题将任务规划各环节与多类型故障注入纳入统一的验证框架中，有助于揭示水域载具在异常扰动下的“规划—执行—重规划”闭环响应规律，为提升自主作业算法的鲁棒性提供可量化的测试手段。

在\textbf{实际应用价值}方面，本课题研发的系统可作为应急处置训练与自主作业能力评估的基础设施。通过可编排的故障注入机制，能够在低成本、无风险的虚拟环境中复现现实中的典型故障，为操作人员提供沉浸式的应急处置演练平台，从而有效弥合算法仿真与实船可靠运行之间的差距。

\section{国内外研究现状}

\subsection{国内研究现状}

在国内研究方面，随着无人船艇技术的快速发展，相关研究正从单一算法向系统级应用与测试评价体系延伸。在任务规划与多艇协同方面，Yang 等人（2025）在《中国舰船研究》提出了一种基于大语言模型驱动的 USV 集群区域覆盖实时决策方法，并在仿真岛屿巡逻场景中验证了多艇协同覆盖与异常事件（如随机 USV 故障）下的任务成功率\cite{yang2025}。该工作表明国内在覆盖搜索与动态重规划算法层面已有面向异常事件的初步设计，但其验证目标仍是“算法有效性”，而非面向应急处置流程的系统级验证。

在技术路线与评价体系方面，吕世兴等（2026）在《中国舰船研究》发表的“无人船艇自主航行决策控制技术综述与展望”中指出，面向开放水域的自主航行技术已日趋成熟，但复杂水域与复杂任务的核心技术仍存在瓶颈，并明确建议构建面向真实场景的标准化仿真和实船测试评价体系\cite{lv2026}。这一观点直接反映了国内在系统级验证与测试评价环境方面的迫切需求，也支撑了本课题构建任务规划验证系统的必要性。

\subsection{国外研究现状}

在国外研究方面，水域载具的任务规划算法与验证技术起步较早，已形成较为完整的理论与仿真体系。在路径规划与算法验证层面，Hashali 等人（2024）系统梳理了无人水面艇（USV）的路径规划算法，指出当前大量工作依赖理想数值仿真验证，并建议未来研究应同时纳入仿真评估与真实水域试验，以全面评价算法的可靠性\cite{hashali2024}。在推进器故障容忍方面，Lagattu 等人（2024）基于深度强化学习研究了 UUV 的推进器故障恢复策略，将故障类型归纳为零输出、常值偏置、时变偏置等，并通过仿真中的故障注入来训练恢复策略\cite{lagattu2024drl}。发表于 \textit{Ocean Engineering}（2023）的 UUV 容错控制综述则系统回顾了执行器故障、传感器故障与通信链路故障的容错方法，将传感器容错策略分为物理冗余与解析冗余两条路径\cite{ftc2023}。这些研究表明，国外在任务规划算法与单一故障容忍控制层面已有深厚积累，但故障多作为控制器设计的约束条件，而非独立的训练变量。

在高保真仿真环境开发方面，国外已推出多款基于虚幻引擎的开源模拟器。HoloOcean 基于虚幻引擎 4（UE4）构建，支持多代理协同、DVL/IMU/声学与光学调制解调器等传感器模型，并实现了基于八叉树的成像声呐模型\cite{holocean}；UNav-Sim 则基于虚幻引擎 5（UE5）提供高保真渲染与 ROS 支持，并包含配套的水下导航数据集\cite{amer2023unav, alvarez2023mimir}。这些高保真模拟器为本课题的系统开发提供了良好的底层运行环境，但其故障注入能力通常仅表现为传感器噪声模型或通信链路丢包，尚未形成面向多类异常的统一编排与应急处置验证框架。

\subsection{研究述评}

综合国内外研究现状可以看出，目前该领域呈现出两条相对分离的线索：任务规划算法研究关注“如何规划得更好”，故障容错研究关注“控制器如何扛住故障”。

具体而言，现有研究存在以下不足：\\
（1）任务规划验证多停留在无故障或单一扰动场景，缺乏对“规划—执行—异常—重规划/返航”闭环的系统化测试；\\
（2）无论是国内还是国外，故障注入研究多以控制器验证为导向，故障类型与注入时机通常服务于特定控制算法的需求，而非面向应急处置流程的独立训练变量；\\
（3）现有基于虚幻引擎的水下模拟器在传感器与通信建模上已具备较好基础，但将多类故障（推进器故障、定位漂移、传感器异常、通信中断）的编排能力与任务规划验证联动的工作尚不充分。

本课题正是针对上述缺口，依托高保真仿真环境，构建一个同时覆盖任务规划验证与多类型故障注入的水域载具仿真验证系统。

\section{研究内容与创新点}

本课题围绕水域载具任务规划验证与故障注入系统展开，研究内容包括以下两个相互支撑的部分：

（1）\textbf{任务规划验证模块}：面向航线规划、覆盖搜索、避障、动态重规划与返航策略，构建可配置的仿真验证流程，支撑自主作业能力的闭环测试。该模块需能够加载任务场景、驱载具执行规划指令，并记录任务执行过程中的轨迹、状态与任务完成情况。

（2）\textbf{故障注入模块}：设计可编排的故障注入机制，支持推进器故障（如推力损失、输出卡死）、定位漂移、传感器异常（如噪声异常、输出冻结）与通信中断四类典型异常的注入，并驱动载具在故障条件下的应急处置流程验证。该模块需支持故障类型、注入时机与持续时间的配置，使故障成为可重复的实验变量。

（创新点待系统实现后补充）

\section{论文组织架构}

本文共分为五章，各章内容安排如下：

第1章为绪论，阐述研究背景与意义，综述国内外在水域载具任务规划与故障注入仿真方面的研究现状，明确研究内容与论文结构。

第2章为系统需求分析与总体设计，分析任务规划验证与故障注入的功能需求与性能指标，给出系统总体架构与模块划分。

第3章为任务规划验证与故障注入模块设计与实现，详细阐述航线规划、覆盖搜索、避障、重规划与返航策略在仿真环境中的验证流程，以及推进器故障、定位漂移、传感器异常与通信中断的注入机制和与任务规划模块的联动方式。

第4章为系统集成与实验验证，介绍系统集成方案，并通过典型任务场景与故障组合实验验证系统的功能完整性与可用性。

第5章为总结与展望，总结本文工作，分析不足并指出后续研究方向。

